\section{Discussion and Threats to Validity}
\label{sec:discussion}

\subsection{A Reversed Finding, Disclosed}

An earlier stage of this investigation, evaluated on a smaller
(121-group), more loosely filtered split, found the trained ranking probe
outperforming zero-shot self-rank by 13 points. That result did not
survive evaluation on the larger, more stringently filtered 1,229-group
pool reported in Section~\ref{sec:results:ranking}, where zero-shot wins
on 6 of 7 rows. We report both results and the reversal explicitly rather
than reporting only the final number, since a reversal under closer
scrutiny is itself informative: it indicates the earlier comparison was
underpowered or measuring an artifact of the looser filter, and a reader
relying on the earlier number alone would draw the wrong practical
conclusion about whether to deploy the probe in place of a zero-shot
baseline.

\subsection{Sample Size and Statistical Power}

Several of the correctness results in Section~\ref{sec:results:correctness}
are drawn from modest sample sizes -- most notably Orchid-BCB at
$n{=}21$ for the hit1-vs-pass-rate comparison. Its $+21.5$pp round-1 gap is
suggestive (Mann-Whitney $p = 0.057$) but not conclusive, and a swing of
one or two cases would move it substantially; we report the wide bootstrap
confidence interval alongside the point estimate for this reason rather
than treating a directionally-interesting result as a settled one. The
headline all-Orchid-combined result ($n{=}213$) is comparatively better
powered and reaches significance by two of three tests applied, but we
regard replication at larger $n$ as necessary before treating the specific
magnitude of the effect as established.

\subsection{Cross-Source Generalization}

The ranking probe's leave-one-source-out performance varies substantially
by held-out source (Table~\ref{tab:ranking}), transferring worst to
Orchid-BCB-Expand. We traced this to the injected ambiguous edit
constituting a small fraction of a long, free-form prompt for
BigCodeBench-derived sources (median edit fraction $2.1$--$2.3\%$ of
prompt words) compared to Orchid-HEval's shorter, more directly-edited
prompts (median $7.1\%$) -- a single pooled or last-token hidden-state
vector has proportionally less of its representational capacity available
to encode a small, diluted edit within a long prompt. This is consistent
with zero-shot self-rank, which re-reads the full prompt text at inference
time rather than compressing it into one fixed-size vector, not suffering
the same degradation on exactly this source. We consider this an
important limitation of single-vector hidden-state probing for long,
free-form specifications, rather than a limitation specific to our probe
architecture, and a natural target for future work (e.g., multi-token or
attention-weighted pooling that preserves more positional information about
where in a long prompt the disambiguating content lives).

\subsection{Choice of Probe-Host Model}

We use Qwen3-8B as the probe-host model throughout due to compute
constraints available to this study, not because it is necessarily the
best choice. A larger probe-host model might encode ambiguity and
candidate-quality information more richly, at the cost of increasing the
per-forward-pass cost that is central to this approach's efficiency
argument; conversely, a smaller model might preserve most of the cost
advantage while sacrificing some signal quality. We do not have data
characterizing this tradeoff and flag it as a direct avenue for future
work rather than assume our specific model choice generalizes.

\subsection{Metric Sensitivity}

Section~\ref{sec:results:correctness}'s central methodological finding --
that the fractional test-case pass rate recovers a significant
ranking-quality effect that the all-or-nothing pass metric does not --
depends on test suites that assign partial credit meaningfully, i.e.,
problems with multiple, at least partially independent test assertions.
For problems with very few test cases, the two metrics converge and this
distinction has little practical effect; we would expect the gap between
metrics to shrink correspondingly on benchmarks with coarser test suites.

\subsection{Absence of a Direct Repeated-Sampling Baseline}

We compare against zero-shot self-rank, which requires a single generation
call per problem, but we do not implement or benchmark a repeated-sampling,
self-consistency-based baseline in the style of ClarifyGPT's underlying
uncertainty-estimation mechanism~\cite{clarifygpt} directly against our
pipeline on the same problems. Any comparison to that class of method in
this paper is therefore a related-work comparison grounded in its published
methodology, not an empirical head-to-head measurement, and should be read
accordingly; we regard implementing and directly benchmarking such a
baseline as important future work to substantiate the cost-efficiency
argument beyond the zero-shot-vs-probe comparison reported here.

\subsection{Construct Validity of the Ambiguity Injections}

Both ClarifyCodeBench and the Orchid-derived sources construct ambiguity by
programmatically editing an originally unambiguous prompt, rather than
sourcing naturally occurring ambiguous specifications. This is a common and
practical approach to constructing matched clear/ambiguous pairs with known
ground truth, but it means our results characterize how well this pipeline
detects and resolves a specific, constructed style of ambiguity, and may
not transfer identically to naturally occurring specification gaps in
real development workflows.
